%!TEX root = ../main_thesis.tex

\chapter{Conclusioni e sviluppi futuri}
\label{cap:conclusioni}

\section{Riepilogo dei contributi}
\label{sec:riepilogo}

Questa tesi ha presentato ROOC, un linguaggio per la modellazione di problemi di
ottimizzazione e l'ambiente che lo accompagna. Il lavoro si è concentrato sul
\emph{motore}: un compilatore che traduce un modello, scritto in una notazione
dichiarativa vicina alla matematica, in una forma lineare risolvibile da un
solver.

I contributi principali possono essere riassunti come segue. È stato progettato
un \emph{linguaggio} dotato di un sistema di tipi e di costrutti di alto livello,
come gli iteratori, le \emph{block function} e i tipi di dato strutturati
(Capitolo~\ref{cap:linguaggio}). È stato realizzato un \emph{compilatore
multi-fase} che, attraverso una sequenza di trasformazioni
(Capitoli~\ref{cap:parsing}--\ref{cap:linearizzazione}), riduce il modello alla
forma standard, e la cui architettura a \emph{pipe}
(Capitolo~\ref{cap:architettura}) rende ogni fase componibile e ispezionabile. È
stato esteso il solver \texttt{microlp}, a partire da un crate esistente
limitato al simplesso, fino a coprire la programmazione lineare intera mista
(Capitolo~\ref{cap:solver}). Infine, il motore è stato reso disponibile come
libreria riutilizzabile e come piattaforma web didattica
(Capitolo~\ref{cap:piattaforma}).

\section{Limitazioni}
\label{sec:limitazioni}

Il lavoro presenta alcune limitazioni note. Sul piano del linguaggio, mancano
operatori logici che permetterebbero di esprimere in modo naturale vincoli di
natura logica. Sul piano della risoluzione, il simplesso implementato nel
progetto è orientato alla chiarezza didattica più che alle prestazioni, e la sua
applicazione diretta è limitata ai modelli di dimensioni contenute; per i
problemi più impegnativi ci si affida ai solver esterni. Infine, l'intera
pipeline è oggi orientata ai modelli lineari, pur essendo l'architettura
predisposta per ospitare in futuro forme diverse.

\section{Sviluppi futuri}
\label{sec:sviluppi-futuri}

Diversi sono i possibili sviluppi. Il più immediato è l'introduzione di operatori
logici, che amplierebbero la classe di problemi di ottimizzazione combinatoria
modellabili in modo diretto. L'architettura a pipe si presta inoltre
all'aggiunta di nuovi solver e di nuove fasi di trasformazione, anche verso
forme non lineari. Sul fronte della piattaforma, l'ulteriore arricchimento degli
strumenti di visualizzazione e della documentazione interattiva ne rafforzerebbe
la vocazione didattica. ROOC si conferma così non solo come uno strumento per
risolvere modelli di ottimizzazione, ma anche come una base estensibile su cui
continuare a costruire.
